\documentclass[../main.tex]{subfiles}
\begin{document}

\section{Bối cảnh của bài toán}

Trong vài năm gần đây, sự phát triển của các mô hình ngôn ngữ lớn (LLM) đã làm thay đổi đáng kể cách con người tương tác với hệ thống trí tuệ nhân tạo. Một trợ lý AI ngày nay không còn chỉ làm nhiệm vụ trả lời câu hỏi ngắn, mà còn được kỳ vọng hỗ trợ lập kế hoạch, đồng hành trong việc học tập, tư vấn mua sắm, theo dõi công việc và duy trì mối quan hệ với người dùng trong thời gian dài. 

Kỳ vọng đó làm xuất hiện một yêu cầu cốt lõi: hệ thống phải nhớ được các thông tin quan trọng về người dùng qua nhiều phiên hội thoại, chứ không chỉ trong phạm vi một phiên ngắn hạn, riêng lẻ.

Trong thực tế, phần lớn các tác nhân hội thoại hiện tại vẫn phụ thuộc rất nhiều vào ngữ cảnh ở ngay phiên hiện tại. Khi lịch sử trò chuyện kéo dài qua nhiều phiên, hệ thống buộc phải cắt bớt nội dung cũ, tóm tắt lại, hoặc chỉ giữ những đoạn được đánh giá là quan trọng nhất. Những cách làm này tuy giúp mô hình tiếp tục vận hành, nhưng đồng thời cũng làm mất dần những chi tiết tưởng như nhỏ mà lại quan trọng với người dùng: một kế hoạch còn dang dở, một thói quen lặp lại, một lần thay đổi quyết định, hay một mối quan hệ nhân quả từng được nhắc thoáng qua trong một phiên từ rất lâu.

\section{Vấn đề cần giải quyết}

\subsection{Bộ nhớ dài hạn của các tác nhân hội thoại không chỉ là lưu lịch sử}

Thoạt nhìn, ta có thể nghĩ rằng bộ nhớ dài hạn chỉ đơn giản là lưu toàn bộ các đoạn chat cũ và tìm lại khi cần. Tuy nhiên, với các hội thoại kéo dài hàng chục phiên hoặc hàng trăm nghìn từ, cách làm đó nhanh chóng bộc lộ hạn chế. Một truy vấn tưởng như đơn giản, chẳng hạn ``người dùng đã từng định làm gì nhưng chưa làm?'', có thể đòi hỏi hệ thống phải nối một sự thật ở phiên đầu, một lần xác nhận ở phiên giữa, và một chi tiết phủ định ở phiên cuối. Nếu chỉ dựa vào truy vấn theo mức độ giống nhau về từ ngữ và ngữ nghĩa, hệ thống rất dễ bỏ lỡ mối liên hệ này.

Khó khăn còn lớn hơn khi câu hỏi không chỉ yêu cầu nhắc lại một chi tiết đơn lẻ mà cần tái dựng một trạng thái nhận thức. Ví dụ, phát biểu ``tôi đang định học Rust'' khác về bản chất với ``tôi đã học Rust'', dù hai câu này có thể cùng nhắc đến một chủ đề. Tương tự, ``người dùng thường làm gì trước cuộc họp'' khác với ``người dùng đã làm gì trong một buổi sáng cụ thể'', và ``làm cách nào để khắc phục lỗi này'' khác với ``lỗi này đã từng xảy ra''. Nếu mọi nội dung đều bị nén về cùng một kiểu biểu diễn duy nhất, sự khác biệt nói trên sẽ bị làm mờ đi.

\subsection{Giới hạn của các giải pháp phổ biến}

Hai hướng tiếp cận đang được dùng phổ biến là lưu trữ văn bản trong kho lưu trữ (thường ở dưới dạng các vector), hoặc để mô hình trích xuất các sự thực (fact) từ hội thoại rồi lập chỉ mục trên các fact đó. Cả hai hướng đều có giá trị thực tế, nhưng chúng cũng có những điểm yếu dễ nhận thấy:

\begin{itemize}
    \item \textbf{Mất mát chi tiết do nén}: Khi hội thoại bị rút gọn thành fact ngắn, nhiều thông tin nguyên bản như con số, tên riêng, sắc thái cảm xúc hoặc ngữ cảnh trích dẫn rất khó được phục hồi đầy đủ.
    \item \textbf{Biểu diễn chưa đủ chuyên biệt}: Các fact về tri thức khách quan, trải nghiệm cá nhân, thói quen, kế hoạch và quy luật nhân quả thường bị đặt chung vào một không gian biểu diễn, dẫn tới truy vấn chưa đúng trọng tâm.
    \item \textbf{Truy vấn đa bước yếu}: Những câu hỏi cần nối nhiều dấu vết nhỏ qua nhiều phiên thường không được xử lý tốt nếu hệ thống chỉ giữ một đường suy luận qua với các bằng chứng mạnh nhất.
    \item \textbf{Định tuyến truy vấn}: Một câu hỏi hỏi ``khi nào'' cần logic truy vấn khác với câu hỏi hỏi ``vì sao'' hoặc ``thường làm gì''. Dùng cùng một cách truy vấn cho mọi loại làm giảm hiệu quả của từng kênh.
\end{itemize}

Những hạn chế trên cho thấy vấn đề không chỉ nằm ở việc mô hình có ``nhớ nhiều hơn'' hay không, mà nằm ở chỗ hệ thống có tổ chức trí nhớ đúng cách hay không. Một bộ nhớ dài hạn hữu ích cần vừa giữ được bằng chứng gốc khi cần, vừa cho phép biểu diễn sự khác nhau giữa tri thức, kinh nghiệm, kế hoạch, thói quen và quan hệ nhân quả.

\section{Mục tiêu và hướng tiếp cận của đồ án}

Từ các quan sát trên, đồ án này đề xuất \textbf{CogMem}, một khung bộ nhớ dài hạn cho tác nhân hội thoại được thiết kế theo hướng kết hợp \textit{khoa học nhận thức} với \textit{đồ thị tri thức}. Ý tưởng xuất phát là: bộ nhớ người không phải một kho lưu trữ đồng nhất, mà là tập hợp của nhiều hệ lưu trữ con cùng tồn tại, mỗi hệ lưu trữ con phù hợp với một loại câu hỏi và một loại suy luận khác nhau. Nếu muốn xây dựng một bộ nhớ dài hạn đáng tin cậy cho tác nhân hội thoại, tác nhân AI hay trợ lý cá nhân, hệ thống cũng cần một cấu trúc biểu diễn phân hóa tương tự.

Theo hướng đó, CogMem theo đuổi bốn đóng góp chính:

\begin{enumerate}
    \item \textbf{Đồ thị bộ nhớ dựa trên khoa học nhận thức}: Thay vì chỉ lưu một loại fact chung, CogMem phân chia trí nhớ thành các nhóm biểu diễn như tri thức khách quan, sự kiện cá nhân, quan điểm, thói quen, ý định tương lai và quan hệ hành động -kết quả.
    \item \textbf{Lưu song song biểu diễn nén và hội thoại gốc}: Mỗi đơn vị bộ nhớ không chỉ chứa fact đã rút gọn mà còn giữ lại một đoạn trích nguyên văn để hỗ trợ trả lời có căn cứ.
    \item \textbf{Lan truyền tín hiệu theo hướng cộng dồn}: Thay vì chỉ chọn một đường có tín hiệu mạnh nhất trong đồ thị, hệ thống cho phép nhiều đường trích dẫn cùng đóng góp để tăng khả năng truy vấn nhiều bước.
    \item \textbf{Định tuyến truy vấn theo bản chất câu hỏi}: Cùng là truy vấn bộ nhớ, nhưng câu hỏi về thời gian, sở thích, nguyên nhân hay kế hoạch tương lai cần được ưu tiên những loại bằng chứng khác nhau.
\end{enumerate}

Cần nhấn mạnh rằng đồ án không giả định các thành phần mới là một sự bảo đảm cho việc mọi bộ dữ liệu chuẩn đều tăng điểm đồng đều. Toàn bộ báo cáo được trình bày theo tinh thần thực nghiệm cẩn thận: kiến trúc được đề xuất vì nó hợp lý về mặt biểu diễn và truy vấn, còn mức độ hiệu quả thực tế của từng thành phần sẽ được kiểm chứng riêng rẽ trên các bộ dữ liệu chuẩn và trên bộ dữ liệu tự xây dựng của dự án.

\section{Phạm vi của báo cáo}

Báo cáo tập trung vào bài toán bộ nhớ dài hạn cho \textbf{tác nhân hội thoại định hướng cá nhân hóa}. Phạm vi này bao gồm:

\begin{itemize}
    \item lưu trữ và truy hồi thông tin người dùng qua nhiều phiên;
    \item hỗ trợ câu hỏi về sự kiện đã xảy ra, sở thích, thói quen, kế hoạch tương lai và quan hệ nhân quả;
    \item đánh giá trên các bộ dữ liệu hội thoại dài tiêu biểu và trên một benchmark phục vụ dự án.
\end{itemize}

Những nội dung ngoài phạm vi chính của báo cáo gồm: tối ưu hóa hạ tầng phân tán ở quy mô sản phẩm hoàn chỉnh, huấn luyện lại mô hình nền, hoặc phát triển đầy đủ một hệ thống điều phối tác vụ phức tạp ngoài bài toán bộ nhớ. Dù vậy, trong nhiều đoạn phân tích, báo cáo vẫn tham chiếu đến các bối cảnh tác nhân AI thực tế để làm rõ vì sao các cấu trúc được đề xuất có ý nghĩa.

\section{Cấu trúc của báo cáo}

Báo cáo được tổ chức lại thành năm chương chính và các phụ lục như sau:

\begin{itemize}
    \item \textbf{Chương 1 -- Đặt vấn đề}: trình bày bối cảnh, bài toán, động lực nghiên cứu và phạm vi của đồ án.
    \item \textbf{Chương 2 -- Cơ sở lý thuyết và các nghiên cứu liên quan}: tổng hợp kiến thức nền tảng về mô hình ngôn ngữ lớn, tác nhân AI, bộ nhớ ngoài cho tác nhân hội thoại, khoa học nhận thức, đồ thị tri thức, các hệ thống bộ nhớ liên quan và hai bộ dữ liệu đánh giá chính.
    \item \textbf{Chương 3 -- Đề xuất CogMem: Khung bộ nhớ dài hạn dựa trên khoa học nhận thức}: mô tả đầy đủ kiến trúc CogMem, luồng xử lý, các lựa chọn biểu diễn, cơ sở khoa học của từng thành phần và các cơ chế truy vấn chính.
    \item \textbf{Chương 4 -- Kết quả thực nghiệm và đánh giá}: trình bày thiết kế thực nghiệm, cách chắt lọc dữ liệu, các thước đo, các cấu hình so sánh, kết quả trên LongMemEval và LoCoMo, cùng phần đánh giá định tính trên CogMem Bench.
    \item \textbf{Chương 5 -- Kết luận và hướng phát triển}: tổng kết đóng góp, những giới hạn hiện tại và các hướng mở rộng thực tế của nghiên cứu trong tương lai.
\end{itemize}

\vspace{1em}
\noindent\textit{Sau khi làm rõ bài toán và mục tiêu của đồ án, Chương 2 sẽ đi vào nền tảng lý thuyết và các nghiên cứu liên quan để xác định chính xác khoảng trống mà CogMem sẽ xử lý.}

\end{document}
